\documentclass[11pt]{article}
\usepackage[margin=0.85in]{geometry}
\usepackage{amsmath,amssymb,booktabs,xurl,hyperref}
\usepackage[T1]{fontenc}
\setlength{\emergencystretch}{2em}
\setlength{\tabcolsep}{4pt}
\title{A Control Lyapunov Function for Recovery Without Collision Risk}
\author{Pan Dynamics Collision Avoidance Research}
\date{October 1, 2026}
\begin{document}
\maketitle

\section{Requirement and result}
The requirement is a control Lyapunov function (CLF) whose decrease condition,
met at every frame without collision risk, dissipates the ego state to the
nominal behaviour: the given path at the reference speed. This needs two
properties.
\begin{enumerate}
\item \emph{Existence:} at every state the ego can be in, some admissible input
satisfies the condition.
\item \emph{Use:} the controller actually issues such an input.
\end{enumerate}
The previous lane quadratic failed the first property outside a small set
($V\le44$ at 8~m/s, $V\le27$ at 15~m/s;
\path{CLF_NO_RISK_DISSIPATION_20261001.tex}). The controller also failed the
second property near the path, through the trust-region trap.

The new CLF is the cost-to-go of a path-guidance feedback $\kappa$. Along
$\kappa$, $V$ falls by exactly the stage cost in the controller's model, so
$\kappa(x)$ is the required input wherever $\kappa$ converges.
\begin{itemize}
\item $\kappa$ converges from all 10{,}800 sampled starts: offsets up to
150~m, every heading, speed factors 0.5--1.15, lateral velocities to 1.5~m/s
and yaw rates to 0.5~rad/s, on straight and 200-m-radius paths at 8 and
15~m/s.
\item It also converges from all 1{,}592 recorded no-risk states of the
previous campaign.
\item Without a target in range, the controller anchors on $\kappa$ and issues
it.
\end{itemize}
In the fourteen-case campaign, every ego that leaves the encounter range now
returns to nominal cruise: 12 of 14 cases recover, at 14.15--46.20~s,
against 3 of 14, at 43.85--62.65~s, before.
\begin{itemize}
\item At all 3{,}452 target-free frames the CLF slack is zero.
\item Between consecutive target-free frames, the CLF measured at the ODE45
plant fell by at least the required amount every time (3{,}433 of 3{,}433).
\item The two crossing cases never leave the encounter range, so this CLF never
applies to them.
\item Encounter-phase behaviour is unchanged, and so are the minimum gaps,
except one case whose first two frames start beyond the range.
\end{itemize}

\section{Construction}
\subsection{The CLF}
Let $f$ be the controller's hold model (one RK4 step of the Fiala bicycle per
50-ms hold) and $\kappa$ the recovery feedback below. Write
$F(x)=f(x,\kappa(x))$ and $e(x)=[e_y,e_\psi,e_v,v_y-v_y^*,r-r^*]$, the
existing path-frame error. The stage cost is $l(e)=e^TQe$ with
$Q=\mathrm{diag}(\mathrm{scales})^{-2}$ from \texttt{cfg.clf}. Define
\[
V(x)=\sum_{k=0}^{K(x)-1}l\bigl(e(x_k)\bigr)+V_f\bigl(e(x_{K(x)})\bigr),\qquad x_0=x,\ x_{k+1}=F(x_k),
\]
where $K(x)$ is the first step with $V_f\le10^{-3}$, and
$V_f(e)=e^TP_fe$ solves $A_{\rm cl}^TP_fA_{\rm cl}-P_f+Q=0$ for the closed-loop
Jacobian $A_{\rm cl}$ at the trim.
\begin{itemize}
\item \emph{Decrease.} $F(x)$ lies on the same closed loop one hold later and
reaches the stop level at the same absolute step, so
\[
V(F(x))=V(x)-l(e(x)).
\]
At every state from which $\kappa$ converges, $\kappa(x)$ therefore decreases
$V$ by the full stage cost. That is the CLF existence property, with region
equal to the region of convergence of $\kappa$.
\item \emph{Positivity.} $V\ge0$, and $V=0$ only at the trim, which is an
equilibrium of $\kappa$.
\item \emph{Dissipation.} If every frame meets
$V(x_{k+1})\le V(x_k)-\eta\,l(e(x_k))$ with $\eta>0$, summing gives
$\eta\sum_k l(e(x_k))\le V(x_0)$. Hence $l(e(x_k))\to0$: the ego converges to
the nominal behaviour, without any fixed exponential rate.
\end{itemize}

\subsection{The recovery feedback}
\texttt{nonlinearBicycleModel.recoveryInput} cascades four loops.
\begin{enumerate}
\item \emph{Guidance.} The course relative to the path is
$\chi=\psi+\arctan(v_y/v_x)-\psi_{\rm path}$, and its desired value is
$\chi_d=-\arctan(e_y/D)$. With exact course tracking,
$\dot e_y=-Ve_y/\sqrt{e_y^2+D^2}$, so $e_y^2$ decreases at every offset. The
approach angle never exceeds $90^\circ$, unlike the lane quadratic's preferred
heading.
\item \emph{Course loop.} The desired yaw rate is
$r_d=\kappa_pV\cos\chi/(1-\kappa_pe_y)+\dot\chi_d-k_\chi\,\mathrm{wrap}(\chi-\chi_d)$.
It is limited to $0.75\mu g/V$, so lateral acceleration stays within
friction. A reversed course turns at the limit.
\item \emph{Yaw-rate loop.} The required front lateral force is
$\bigl(I_zk_r(r_d-r)+l_rF_{yr}\bigr)/l_f$, capped at 90\% of the
combined-slip capacity. The front slip angle comes from the exact Fiala
inverse, $|F|/F_{\max}=1-(1-s)^3$ with $\tan|\alpha|=3F_{\max}s/C$, so the tire
never passes its force peak.
\item \emph{Speed loop.} The braking ratio is $b=b^*+k_v(v_x^*-v_x)$, limited to
$\pm0.35$.
\end{enumerate}
The inverse neglects the front longitudinal force and $\cos\delta$. A constant
steering offset, computed once per configuration and curvature, therefore makes
$\kappa(\text{trim})=u^*$ exactly. The offset is 0 on straight paths, and
$-1.68\times10^{-5}$~rad (8~m/s) and $-2.68\times10^{-5}$~rad (15~m/s) on the
200-m curve.
\begin{center}\footnotesize
\begin{tabular}{llll}
\toprule
Parameter (\texttt{cfg.recovery}) & Value & Parameter & Value\\
\midrule
lookahead $D=\max(8,\,1.5\,v_{\rm ref})$ & 12~m / 22.5~m & speedGain $k_v$ & 0.5 per m/s\\
courseGain $k_\chi$ & 1.5~s$^{-1}$ & brakingRatioLimit & 0.35\\
yawRateGain $k_r$ & 10~s$^{-1}$ & decreaseFraction $\eta$ & 0.5\\
lateralAccelerationFraction & 0.75 & stopValue & $10^{-3}$\\
frontForceFraction & 0.9 & maximumSeconds & 120~s\\
\bottomrule
\end{tabular}\end{center}

\section{Sampled certificate}
\path{scripts/certifyRecoveryClf.m} evaluates $\kappa$ in the controller's
model.
\begin{center}\footnotesize
\begin{tabular}{rrrrrrrrr}
\toprule
Speed & Curvature & Starts & Pass & Entry median / max (s) & Max $|a_y|$ & Min $v_x$ & Spectral radius & Lyapunov residual\\
\midrule
8 & 0 & 2700 & 2700 & 12.65 / 30.05 & 7.05 & 4.13 & 0.9682 & $2.8\times10^{-15}$\\
8 & 0.005 & 2700 & 2700 & 12.65 / 30.15 & 7.05 & 4.13 & 0.9682 & $1.0\times10^{-15}$\\
15 & 0 & 2700 & 2700 & 15.80 / 26.95 & 7.00 & 7.62 & 0.9694 & $1.5\times10^{-15}$\\
15 & 0.005 & 2700 & 2700 & 15.75 / 27.05 & 7.00 & 7.62 & 0.9695 & $2.6\times10^{-15}$\\
\bottomrule
\end{tabular}\end{center}
\begin{itemize}
\item \emph{Grid.} The starts cover lateral offsets $\{0,\pm1,\pm5,\pm15,\pm30,\pm60,\pm100,\pm150\}$~m and headings
every $20^\circ$ relative to the path. Ten combinations of speed factor, lateral
velocity and yaw rate are crossed with them.
\item \emph{Pass criterion.} A start passes when all five errors stay inside the recovery
tolerances (0.1~m, $1^\circ$, 0.1~m/s, 0.05~m/s, 0.01~rad/s) for 5~s. ``Entry''
is the time the errors first reach them. Acceleration is in m/s$^2$ and speed in m/s.
\item \emph{Results.} No start leaves the tire or model domain. The peak yaw rate is
1.17~rad/s and the peak lateral velocity 1.51~m/s.
\item \emph{Identity check.} On every tenth start (1{,}080 starts),
$V(F(x))=V(x)-l(e(x))$ holds to a relative error of at most
$9.5\times10^{-16}$.
\item \emph{Recorded states.} From every tenth no-risk frame of the 5-cm
baseline replays (1{,}592 states, $e_y$ from $-177$ to $533$~m), $\kappa$
enters and holds the tolerances within 45.05~s, with no domain failure.
\end{itemize}

\section{Use in the controller}
When the measured state is beyond the 30-m encounter range, or there is no
target, the solve has no collision rows. Three things change.
\begin{itemize}
\item \emph{Anchor.} The linearization anchor is \texttt{recoveryFeedbackRollout}:
$\kappa$ for the primary horizon, then the existing terminal completion. Its
first input is $\kappa(x_0)$. The shifted previous plan, which caused the trap,
is not used.
\item \emph{CLF stage.} The stage requires
$V(x_1)\le V(x_0)-\eta\,l(e(x_0))+\rho$.
\begin{itemize}
\item $V(x_0)$ comes from one rollout.
\item Near the anchor's first node, $V$ is the Gauss--Newton cone $c+|a+R\,\Delta x_1|^2$. Its Jacobian comes from six rollouts of equal length started at perturbed states.
\item At zero correction the cone equals $V(x_0)-l$, so $\rho=0$ is feasible with margin $(1-\eta)\,l$.
\end{itemize}
\item \emph{Third stage.} Subject to both achieved slack levels, a third
problem minimizes the weighted input deviation from the anchor, with weight
$10^3$ on the first input. The issued input is therefore $\kappa(x_0)$ whenever
the other constraints admit it, and $V$ falls by exactly $l$ in the hold
model.
\end{itemize}
The cone alone is not trusted for other inputs. For corrections across the
whole trust box, its error reached twice the stage cost.

While a target is within range, nothing changes: shifted-plan anchor, PCBF
stage, lane-quadratic CLF stage, and no third stage. The continuation state
version is now 64.

\section{Validation}
\paragraph{Tests.} The full MATLAB suite passes, 568 of 568. That includes the eleven
controller test classes (223 tests) and the new
\path{tests/recoveryClfTest.m}, which checks:
\begin{itemize}
\item the exact trim equilibrium;
\item the Lyapunov equation;
\item the cost-to-go identity;
\item convergence from distant and reversed starts;
\item a 21-frame target-free closed loop.
\end{itemize}
The Python audit tests pass, 20 of 20.

Some expectations changed. Target-free frames now report
\texttt{recoveryFeedbackRollout}, three solver calls and the
anchor-deviation stage. The shifted-input, restart and steering-cap tests now
place a target within range, where the two-stage behaviour is unchanged.

\paragraph{Campaign.} The fourteen collision-threat cases of the 5-cm campaign
were rerun with the changed controller. Fixtures, driver, ODE45 plant, 80-s
window and 5-s recovery dwell are unchanged. In the table, ``free'' frames are
target-free frames, and ``meets'' is the share of consecutive free frames
whose CLF, measured at the plant states, fell by at least $\eta\,l$.
\begin{center}\footnotesize
\setlength{\tabcolsep}{3pt}
\begin{tabular}{rlrrrrrrrr}
\toprule
 & & \multicolumn{2}{c}{Recovered at (s)} & \multicolumn{2}{c}{Min gap (m)} & Free & Meets & Max $|\Delta V_{\rm plant}-l|$ & Median free\\
Speed & Scenario & before & now & before & now & frames & $\eta\,l$ & $/\,l$ & frame (s)\\
\midrule
8 & headOn & 57.10 & 18.05 & 0.603 & 0.603 & 264 & 100\% & 0.038 & 0.19\\
8 & acceleratingHeadOn & 62.65 & 17.40 & 0.430 & 0.430 & 266 & 100\% & 0.048 & 0.20\\
8 & brakingLead & 43.85 & 33.45 & 1.062 & 1.062 & 248 & 100\% & 0.004 & 0.17\\
8 & crossing & --- & --- & 0.742 & 0.742 & 0 & --- & --- & ---\\
8 & turningCrossing & --- & 40.40 & 0.440 & 0.440 & 391 & 100\% & 0.038 & 0.20\\
8 & curvedHeadOn & --- & 18.00 & 0.707 & 0.707 & 259 & 100\% & 0.031 & 0.19\\
8 & curvedCrossing & --- & 46.20 & 0.589 & 0.589 & 282 & 100\% & 0.022 & 0.23\\
15 & headOn & --- & 14.50 & 0.035 & 0.035 & 231 & 100\% & 0.0007 & 0.19\\
15 & acceleratingHeadOn & --- & 14.15 & 0.031 & 0.031 & 229 & 100\% & 0.0005 & 0.21\\
15 & brakingLead & --- & 45.25 & 0.710 & 0.710 & 238 & 100\% & 0.0003 & 0.19\\
15 & crossing & --- & --- & 0.409 & 0.409 & 0 & --- & --- & ---\\
15 & turningCrossing & --- & 39.20 & 0.128 & 0.128 & 439 & 100\% & 0.0003 & 0.21\\
15 & curvedHeadOn & --- & 15.45 & 0.854 & 0.844 & 233 & 100\% & 0.0009 & 0.18\\
15 & curvedCrossing & --- & 44.30 & 0.151 & 0.151 & 372 & 100\% & 0.0004 & 0.20\\
\bottomrule
\end{tabular}
\end{center}
Notes on the table:
\begin{itemize}
\item ``Before'' is the 5-cm campaign (\path{COLLISION_MARGIN_5CM_20261001.tex}).
\item ``Max $|\Delta V_{\rm plant}-l|/l$'' compares the CLF decrease measured
between plant states with the hold model's exact decrease $l$.
\end{itemize}

\paragraph{Every ego that leaves the encounter range recovers.} This holds for
all twelve such cases.
\begin{itemize}
\item The head-on family confirms recovery at 14.15--18.05~s, within 1.2--5.1~s
of the earliest possible confirmation at 13~s.
\item Braking lead recovers at 33.45 and 45.25~s.
\item The turning and curved crossing cases recover at 39.20--46.20~s.
\end{itemize}
Nine of these cases never recovered before. The three that did now recover
10--45~s sooner.

\paragraph{The CLF condition held at every target-free frame.} The CLF slack
is zero at all 3{,}452 such frames. Between consecutive target-free frames,
$V$ fell by at least $\eta\,l$ in all 3{,}433 cases, measured at the ODE45 plant
states. The plant decrease differs from the hold model's exact decrease $l$ by at most
4.8\% of $l$ at 8~m/s and 0.09\% at 15~m/s, so the 50\% margin is never
approached. Summed over a run, this is the dissipation argument of Section~2
realized in the plant.

\paragraph{The encounter phase is unchanged.} Minimum gaps equal the
baseline's, except in 15-m/s curvedHeadOn (0.844 against 0.854~m). There the
target is 31.8 and 30.6~m away in the first two frames, which are therefore
target-free and now use the recovery anchor. The flow restarts are also
unchanged: one in each 15-m/s curved case.

\paragraph{The crossing cases are not reached.} Both crossing egos still escort
the target inside the 30-m range for the whole 80~s. They end at $-621.0$ and
$-629.0$~m, so the recovery CLF never applies. That trap is in the encounter
phase, which this change leaves as it was.

\paragraph{Cost.} A target-free frame needs seven feedback rollouts and three
cone solves. Its median controller time is 0.17--0.23~s and its maximum
0.39~s, measured with other MATLAB jobs running.

\section{Limits}
\begin{itemize}
\item The existence property rests on the convergence of $\kappa$, which is
sampled, not proved, on the stated grid and recorded states. Starting speeds
below about 4~m/s and curvature charts with $1-\kappa_pe_y<0.1$ are outside
it.
\item The exact decrease holds in the hold model for frames that issue
$\kappa(x_0)$. The ODE45 plant is measured, not covered by the claim.
\item The anchor-deviation stage keeps the first input on $\kappa$ unless other
constraints prevent it, and such frames are not covered by the identity.
\item Nothing is claimed while a target is within range. Safety keeps priority,
and an ego that never leaves the range never reaches this CLF.
\item Controller timings were measured with other MATLAB jobs on the machine;
they are not benchmarks.
\end{itemize}

\section{Artifacts}
The controller changes are in \path{controller/nonlinearBicycleModel.m},
\path{controller/solvePredictiveControl.m},
\path{controller/collisionAvoidanceController.m} and
\path{config/collisionAvoidanceControllerConfig.m}. The derivation is in
\path{controller/RECOVERY_CLF.md}.
\path{report/RECOVERY_CLF_20261001/} holds:
\begin{itemize}
\item the certificate tables;
\item the recorded-state check;
\item the campaign comparison;
\item methods, logs and a manifest.
\end{itemize}
The campaign JSON/MAT files remain in
\path{/home/zai/.cache/collisionAvoidance/recovery-clf-20261001/}.
\end{document}
